\begin{figure}[htbp]
\centering
\begin{tikzpicture}[
    font=\small,
    caja/.style={draw=#1, fill=#1!8, rounded corners=3pt, align=center,
                 minimum width=3.6cm, minimum height=1.2cm, inner sep=3pt},
    caja/.default=subsectionblue,
    flecha/.style={-{Stealth[length=2mm]}, thick, draw=gray!70!black},
    et/.style={font=\scriptsize, text=gray!55!black, align=center, inner sep=2pt},
]

\node[caja=gray]          (sql)   at (0,0)      {Resultado SQL\\[-1pt]\scriptsize todas las filas};
\node[caja=ugrred]        (llm)   at (5.8,2.1)  {Modelo\\[-1pt]\scriptsize diseña la representación};
\node[caja]               (merge) at (5.8,-2.1) {\texttt{graficos.py}\\[-1pt]\scriptsize inyecta los datos};
\node[caja]               (ui)    at (11.4,-2.1){Interfaz\\[-1pt]\scriptsize Vega-Lite o tabla};

\draw[flecha] (sql.north) |- node[et, above, pos=0.72]
      {1. nombres de columna\\y muestra de 3 filas} (llm.west);
\draw[flecha] (llm.south) -- node[et, right]
      {2. tipo, estructura\\y razonamiento,\\sin ninguna cifra} (merge.north);
\draw[flecha] (sql.south) |- node[et, below, pos=0.72]
      {3. el 100\,\% de las filas} (merge.west);
\draw[flecha] (merge.east) -- node[et, above] {4. WebSocket} (ui.west);

\end{tikzpicture}
\caption{Motor visual generativo. El modelo solo decide cómo representar la
información; los valores los pone el \emph{backend}, que ya los tiene. Esto
ahorra unos 700 \textit{tokens} de generación por elemento visual y elimina de
raíz que el modelo se invente cifras al copiarlas.}
\label{fig:generative_ui}
\end{figure}
